\subsection*{S1.20 Natural SNV annotations and descriptive conservation analyses}

\paragraph{Cohort and measured outcome.} The natural-variant analysis included all 7,262 previously qualified BEAPR events across 140 RBPs and 200 RBP--cell tasks \citep{yang2019asb}. An event combines variant, RBP and cell context; it is not an independent biological replicate. The existing data partition comprises 5,409 training, 1,457 validation and 396 chromosome-8/15 events. The last group contains 287 unique GRCh38 loci/REF--ALT variants, 85 RBPs and 119 RNA-region clusters. Both the complete cohort and the chromosome-8/15 subset had been examined during earlier analyses. The source supplies significant ASB calls (source FDR $<0.1$) and normalized major/minor allele identities, but not usable numerical normalized allele-specific effects. We evaluate preferred direction, not quantitative binding change, pathogenicity or ASB detection. Raw replicate counts provide depth diagnostics and never replace the normalized preferred-allele label. No negative-class AUROC or AP is calculated.

Variant mapping required unique, continuous liftover of the complete window and reciprocal mapping of the central position. Each 101-nt allele pair differs only at zero-based index 50. The genome-plus REF allele belongs to the two experimentally compared alleles in every qualified event. Annotation queries use GRCh38 genome-plus REF/ALT; model inputs remain RNA-strand oriented. Both alleles receive the same reference hg38 PhyloP100way vector. This is fixed reference conservation, not a variant-specific evolutionary score.

\paragraph{Predictions, controls and weighting.} For each existing checkpoint, the predicted contrast is the major-preference minus minor-preference count-head output. Contrasts above $10^{-7}$ receive credit one, absolute contrasts at most $10^{-7}$ receive 0.5, and negative contrasts below $-10^{-7}$ receive zero. This local analysis threshold differs from the historical $10^{-6}$ convention in S1.14 and is recorded in the frozen protocol. Three permuted-PhyloP predictions are averaged before taking their sign; this differs from averaging three direction accuracies. The constant REF rule receives one when the reference allele is the source's normalized preferred allele and zero otherwise. Output differences are model log-total-count responses, not calibrated affinity or measured allele-count ratios. Full and Sequence reuse single checkpoints whose complete training histories are not verified as matched. Flat and permuted PhyloP modify the same Full checkpoint at inference; they are not retrained controls.

The two thresholds $10^{-7}$ and $10^{-6}$ changed zero event credits in the complete 7,262-event prediction archive used here; neither threshold produced a tie in any arm. This threshold comparison applies to the prediction archive analysed here; numerical identity with earlier inference archives was not assumed.

Direction credit is averaged over events within task, tasks within RBP and then RBPs equally. Conservation-quartile intervals retained in Data S25 and Figures~S23/S25, and accuracy intervals in Figure~S22, use 5,000 paired RBP bootstrap draws (seed 20261012). Figure~5d displays the five-arm aggregate with the frozen two-way sensitivity intervals; the paired Full-minus-Sequence contrast uses the same interval scheme. A positive-exponential multiplier assigns weights to RNA-region clusters within tasks, averages tasks equally within RBP, and applies independent RBP weights (2,000 draws, seed 20261015). Data S25 contains both interval sets, while Figure~S25 retains the full quartile matrix. Gene-connected regions need not identify one uniquely resolved mature RNA. The RBP bootstrap alone does not represent all shared-locus dependence.

\paragraph{Conservation association and uncertainty.} Within each cohort, center PhyloP is ranked with average ranks for ties and standardized to mean zero and standard deviation one. The response-magnitude outcome is the similarly standardized rank of the absolute predicted allele-paired output difference. The correctness outcome is the original 0/0.5/1 credit, so its slope is change in credit per standard deviation of ranked PhyloP. Nuisance regressors are task indicators; whole-window and central-seven-nt GC fractions; one-hot bases at the five central positions; genome-plus REF$>$ALT mutation class; RNA strand; log(1+total retained raw reads); and the source's reference-allele preference. Redundant categorical effects are handled by the pseudoinverse. Local base encoding is a limited sequence control, not an experimentally measured cognate-motif score.

Primary slope intervals use first-order two-way positive-exponential multipliers of regression scores, clustered by RNA region and RBP (2,000 draws, seed 20261013). These are approximate exploratory intervals, not 2,000 complete nuisance-model refits. The approximation replaced the planned full-refit procedure to limit computational cost; the analysis manifest records this change. Two hundred actual weighted least-squares refits (seed 20261014) checked the same eight slopes; all chromosome-8/15 intervals crossed zero. A later exploratory sensitivity added NCBI region classes, including unavailable annotation, and performed 200 further weighted refits (seed 20261016). The annotation hierarchy is coding, splice, UTR, noncoding exon, intronic, other annotated, unavailable; it is a nuisance grouping, not measured SNV function. Full's region-adjusted slopes were 0.031 for correctness (95\% interval $[-0.052,0.113]$) and 0.046 for ranked magnitude ($[-0.139,0.191]$). Paired Full-minus-Sequence slope differences also remained compatible with zero.

Center-PhyloP quartiles are defined from unique loci separately in each cohort without outcomes; ties remain together. Unequal event and RBP coverage is retained in Data S25. In the full development-exposed 7,262-event cohort, predicted-magnitude slopes were 0.086 for Full (approximate interval 0.044--0.130) and 0.045 for flat PhyloP (0.004--0.087). This development-set association does not establish improved allele ordering or a conservation-specific effect, and did not reproduce as a stable chromosome-8/15 association. All controls and strata are reported; example selection used the coverage-based rule below.

\paragraph{Exact public annotations and missingness.} Annotation was performed after fixing all 287 unique chromosome-8/15 REF/ALT variants, and annotation availability did not change the 396-event evaluation. Official Ensembl VEP GRCh38 was queried first (software release 116, REST release 15.12). After repeated HTTP500 failures, official NCBI Variation SPDI/RefSNP services supplied exact allele records. SPDI queries use accessions NC\_000008.11 and NC\_000015.10 and zero-based positions. Inclusion requires exact REF/ALT SPDI in the primary GRCh38 chromosome placement. Each population observation is checked for equivalence to that allele against retained placements; a common rs identifier alone is insufficient.

NCBI returned exact matches for 279/287 variants. Eight queries returned HTTP404 twice and remain unavailable. Original 1000 Genomes ALT allele counts and chromosome totals were available for 241 variants/316 ASB events. ALT frequency is their ratio, with source version, numerator and denominator retained. There were 239 biallelic and two triallelic study observations; ALT frequency need not equal minor-allele frequency. Missing values are not zero. Ensembl succeeded for 110 variants and failed for 177; 101 variants had 1000 Genomes frequencies from both services, with maximal absolute difference $4.89\times10^{-5}$, consistent with rounding. Exact NCBI counts are the primary source. Forty-four observations from other population studies lacked directly retained equivalent placements and were excluded only from those annotation records, without dropping any ASB event or primary 1000 Genomes observation.

Thirty-one variants had allele-specific ClinVar links. Direct ClinVar esummary records were matched by canonical SPDI for all 31, retaining VCV version, current classification, review status and evaluation date. Current classifications were Benign (23), Benign/Likely benign (six), Likely benign (one) and Conflicting classifications of pathogenicity (one). Older RefSNP RCV terms remain separately recorded. Lack of a linked entry does not demonstrate absence from current ClinVar. These classifications and population frequencies supply provenance and context; neither is a measured RBP-binding effect.

TCGA mutation records, as described in official GDC documentation, lack the matched allele-specific RBP-binding preferences required for this endpoint. They are excluded from this 396-event direction-accuracy analysis; the separate current whole-cohort annotation join is described in S1.21 and does not create new binding labels. No controlled individual-level data were requested. Cached official replies, URLs, retrieval records and failures are retained in Data S25.

\paragraph{Illustrative natural variant.} The example maximizes distinct-RBP coverage over the full chromosome-8/15 cohort, breaking ties by greater center PhyloP, then chromosome, coordinate and allele. This rule uses coverage/reference information, not model advantage. It selects chr8:1,844,440 C$>$T in the ARHGEF10/KBTBD11-OT1 interval. All ten K562 RBP events appear alphabetically. The source prefers C in all ten; Full recovers four and Sequence six. Center PhyloP is $-3.501$, so the site is not presented as strongly conserved. The example connects a natural variant to measured allelic preference and model response. It is descriptive and does not establish a causal role for conservation.

\paragraph{Display of retained responses.} Current Figure~5d preserves all 396 Full/Sequence source-preferred-minus-other response pairs, without jitter or event filtering. Both axes use the same asinh transform with linear width 0.02, raw limits $[-0.2,1.8]$ and ticks in original log-count units. It separates Full-only agreement (43), Sequence-only agreement (61), both models correct (168) and neither model correct (124), with no ties. These unweighted event counts differ from the RBP-equal accuracy and frozen two-way intervals reported beside them. Earlier joint densities and marginals remain preserved in prior revision source/display files; the complete four-arm event field and conservation summaries remain in Figures~S25/S28 and Data S25. The drawing adds no endpoint, smoothing or significance test.
